\section{The Mechanism: From Global Algebra to Local Consistency}\label{sec:mechanism}\label{sec:setup}

To guarantee the fidelity of a hierarchical pipeline, we must first define what it means for a summary to be ``correct'' in a formal sense. We begin with the formal setup, describe the algebraic ideal we strive for, and then introduce the consistency conditions that effectively force a stochastic LLM to approximate that ideal.

\subsection{Setup: Strings, Oracles, and Metrics}

We treat every document as a finite sequence of tokens. Let $\Strings$ denote the set of such strings, write $u\concat v$ for concatenation, and let $\emptystr$ denote the empty string. The task of interest is represented by a fixed oracle $\fstar:\Strings\to\Y$ that maps strings to task values (labels, structured tuples, embeddings, etc.). In many applications $\Y$ is best viewed as a \emph{state of accumulated evidence}: a running list of actors, claims, or scores that grows as we read more text. We therefore assume (and later audit) that $\fstar$ is well defined on substrings and that $\Y$ carries a neutral element $\emptyset$ representing ``no evidence yet,'' so that leaves with no signal map to $\emptyset$ while informative leaves update the evidence state. The quality of a summary is measured in the oracle space via a (pseudo)metric $\metric:\Y\times\Y\to[0,\infty)$.

\paragraph{A strong locality assumption.}
Treating $\Y$ as an accumulated-evidence state is powerful but brittle: it requires every leaf summary to carry the sufficient statistics needed to evaluate \emph{any} continuation. For narrative tasks the oracle on a substring can be undefined without extra latent variables (e.g., ``Did the bill become law?'' before the court decision is read). We therefore flag this ``oracle on substrings'' premise as an empirical hypothesis rather than an axiom. Section~\ref{sec:audit} spells out the concrete audit we run to stress-test it: we repeatedly probe Condition~\eqref{eq:leaf_compatibility} along realized boundaries and treat persistent violations as evidence that the task cannot be decomposed at the chosen block size. In practice we expect to pair these audits with domain-specific diagnostics (Section~\ref{sec:examples}) and, when necessary, expand the summary format until it truly functions as a sufficient statistic for $\fstar$.

\paragraph{Scope: decomposable tasks.}
Throughout we restrict attention to \emph{decomposable} objectives where the oracle can be evaluated on bounded spans: coding tuples, structured extraction, counting, and guided abstractive summaries whose rubric specifies the state that must propagate. Open-ended ``summarize this document however you like'' goals fall outside our model because no finite $\Y$ suffices. Whenever practitioners apply the framework to guided summaries, they must pin down the rubric (actors, claims, stance, etc.) so that $\fstar$ is literally the function being preserved.

A summarizer $g:\Strings\rightsquigarrow\Strings$ is a (possibly stochastic) map. When $g$ is stochastic, we evaluate its fidelity in expectation. To keep notation minimal, we write $\E[\cdot]$ to average over the internal randomness of $g$ (seeds, temperature). A summary $z$ is considered a valid substitute for document $x$ if they are indistinguishable under the oracle metric,
\[
\metric\big(\fstar(z), \fstar(x)\big)\approx 0.
\]

Downstream learning signals only depend on the oracle. Let $\sstar:\Strings\times\A\to\mathbb{R}$ denote an arbitrary supervision signal (a score, reward, or loss) and let $\A$ be the action or label space. We assume this supervision factors through the task value.

\begin{assumption}[Conditional factorization]\label{ass:CF}
There exists a measurable $\bar Q:\Y\times\A\to\mathbb{R}$ such that for every action $a\in\A$,
\[
\E\!\left[\sstar(X,a)\mid \fstar(X)\right]=\bar Q(\fstar(X),a)\quad\text{almost surely.}
\]
\end{assumption}

This assumption covers standard classification (the loss depends only on the latent label), preference learning (the reward depends on the task semantics), and noisy annotators (the conditioning is on $X$). Once $\fstar(X)$ is known, there is no additional supervision signal left in the raw phrasing.

The next concept lets us reason modulo oracle equivalence. It mirrors the standard ``almost sure equality'' in probability theory but is tailored to the task metric.

\begin{definition}[Oracle equivalence]\label{def:oracle-equivalence}
Two strings $u,v\in\Strings$ are \emph{equivalent at the oracle}, denoted $u\sim v$, if their oracle outputs are indistinguishable:
\[
u\sim v
\;\Longleftrightarrow\;
\metric\big(\fstar(u),\,\fstar(v)\big)=0.
\]
All subsequent conditions are stated modulo this relation. Because $\metric$ is a (pseudo)metric, $\sim$ is reflexive, symmetric, and transitive. Propagating $\sim$ through concatenation requires a context-compatibility property of the oracle:
\end{definition}

\begin{assumption}[Context-compatible oracle]\label{ass:context}
For every context $w\in\Strings$, $u\sim v$ implies both $w\concat u \sim w\concat v$ and $u\concat w\sim v\concat w$.
\end{assumption}

Appendix~\ref{sec:global} records structural conditions on $g$ and $\fstar$ that guarantee Assumption~\ref{ass:context}. We state it here to make the dependency explicit: without context compatibility the equivalence relation would not be a congruence and local checks would fail to propagate. In deployments we treat Assumption~\ref{ass:context} as falsifiable. The merge-consistency checks in Section~\ref{sec:audit} deliberately splice summaries into fresh contexts to detect when oracle values drift under concatenation; persistent failures trigger either larger leaf blocks or richer summary schemas. When $g$ is stochastic, all statements hold almost surely with respect to its internal randomness.

In practice this is exactly the prompt-engineering goal for $g$: when the summarizer keeps all task-relevant content at the leaves (Condition~\eqref{eq:leaf_sufficiency}) and behaves consistently inside the audit (Condition~\eqref{eq:leaf_compatibility}), it should also be a valid substitute when the surrounding context references that content later. The merge-consistency condition formalizes this intuition and lets us turn local checks into guarantees about the raw document.

\begin{definition}[Oracle sufficiency in expectation]\label{def:expect-suff}
The summarizer $g$ is \emph{$\fstar$-sufficient in expectation} if $\E[\metric(\fstar(g(x)),\fstar(x))]=0$ for all $x\in\Strings$.
\end{definition}

\begin{definition}[Summary idempotence]\label{def:resum-stable}
The summarizer $g$ is \emph{summary-idempotent} if every string already on the range of $g$ keeps its oracle under another call to $g$: for any random $Z$ supported on $\operatorname{range}(g)$,
\[
\E\big[\metric(\fstar(g(Z)),\fstar(Z))\,\big|\,Z\big]=0.
\]
\end{definition}

\begin{definition}[Exact classes $\Gstar$ and $\Gstar_{\mathrm{stab}}$]\label{def:exact-suff-stable}
Let $\Gstar:=\{g:\metric(\fstar(g(X)),\fstar(X))=0\ \text{a.s.}\}$ and $\Gstar_{\mathrm{stab}}:=\{g\in\Gstar:\text{$g$ also satisfies Definition~\ref{def:resum-stable}}\}$.
\end{definition}

Doob--Dynkin factorization implies $g\in\Gstar$ iff the oracle can be recovered as a measurable function of $g(X)$. In classical statistics this is exactly the definition of a sufficient statistic: the summary $g(X)$ captures all information about $\fstar(X)$, so conditioning on $g(X)$ leaves no task signal behind. Readers who prefer an intuition-first treatment can consult \citet{doobdynkin-readable}, which emphasizes this sufficiency viewpoint. Stability simply says re-applying $g$ cannot change that oracle. Figure~\ref{fig:gstar} visualizes the equalities enforced by these requirements.

\begin{figure}[t]
    \centering
    \begin{tikzpicture}[node distance=2cm]
    \node[doc] (X) {$X$};
    \node[sum, right=of X] (gX) {$g(X)$};
    \node[sum, right=of gX] (ggX) {$g(g(X))$};
    \node[op, below=of X] (fX) {$\fstar(X)$};
    \node[op, below=of gX] (fgX) {$\fstar(g(X))$};
    \node[op, below=of ggX] (fggX) {$\fstar(g(g(X)))$};
    \draw[->, dotted] (X) -- (gX);
    \draw[->, dotted] (gX) -- (ggX);
    \draw[->] (X) -- (fX);
    \draw[->] (gX) -- (fgX);
    \draw[->] (ggX) -- (fggX);
    \draw[<->, dashed] (fX) -- node[below] {\tiny $\metric$-equal if $g\in\Gstar$} (fgX);
    \draw[<->, dashed] (fgX) -- node[below] {\tiny $\metric$-equal by idempotence} (fggX);
    \end{tikzpicture}
    \caption{Oracle-preserving normal form. Sufficiency forces $\fstar(g(X))$ to equal $\fstar(X)$; on-range stability keeps additional summaries inert.}
    \label{fig:gstar}
\end{figure}

\subsection{Document Decomposition}

Long documents routinely exceed the model's context window, so we cut $x$ into contiguous blocks $(b_1,\ldots,b_k)$ and summarize them along a realized full binary tree $T$. Every leaf stores the raw substring $b_i$ plus its first-pass summary $s_i=g(b_i)$. For any node $u$ we keep $S(u)$ as the latent span of raw text underneath it; we only need this symbol when comparing back to the original document in global arguments. Operationally the model never touches $S(u)$ once we leave the leaves. The only thing that moves up the tree is the stored summary $s_u$. Leaves set $s_u=g(b_i)$, and each internal node $u$ simply reads the strings from its two children, concatenates them, and runs the summarizer again: $s_u := g(s_{u_L}\concat s_{u_R})$. Traversing $T$ bottom-up produces the one-pass output $Z^{(1)}(x):=g(\text{root}(T))$, and optional clean-up passes apply the same rule recursively via $Z^{(R+1)}(x):=g(Z^{(R)}(x))$. This ``summarize, concatenate, summarize again'' recipe is exactly what the audit later inspects.

\begin{figure}[t]
    \centering
    \begin{tikzpicture}[x=1.0cm, y=1.2cm]
  \tikzset{selected/.style={fill=yellow!30, draw=orange!80!black, very thick}}
  \foreach \i [evaluate=\i as \xx using 2*(\i-1)] in {1,...,8}{
    \node[doc] (x\i) at (\xx,0) {$x_{\i}$};
    \node[sum] (g\i) at (\xx,1.4) {$s_{\i}=g(x_{\i})$};
    \draw[->] (x\i) -- (g\i);
  }
  \foreach \a/\b in {1/2,3/4,5/6,7/8}{
    \node[plus] (p\a\b) at ($ (g\a)!0.5!(g\b) + (0,1.2) $) {$\concat$};
    \node[sum]  (s\a\b) at ($ (p\a\b) + (0,1.2) $) {$s_{\a\b}=g\!\big(s_{\a} \concat s_{\b}\big)$};
    \draw[->] (g\a) -- (p\a\b);
    \draw[->] (g\b) -- (p\a\b);
    \draw[->] (p\a\b) -- (s\a\b);
  }
  \foreach \L/\R in {12/34,56/78}{
    \node[plus] (p\L\R) at ($ (s\L)!0.5!(s\R) + (0,1.2) $) {$\concat$};
    \node[sum]  (s\L\R) at ($ (p\L\R) + (0,1.2) $) {$s_{\L\R}=g\big(s_{\L} \concat s_{\R}\big)$};
    \draw[->] (s\L) -- (p\L\R);
    \draw[->] (s\R) -- (p\L\R);
    \draw[->] (p\L\R) -- (s\L\R);
  }
  \node[plus] (pall) at ($ (s1234)!0.5!(s5678) + (0,1.2) $) {$\concat$};
  \node[sum]  (sall) at ($ (pall) + (0,1.2) $) {$s_{\mathrm{root}}=g\big(s_{1234} \concat s_{5678}\big)$};
  \draw[->] (s1234) -- (pall);
  \draw[->] (s5678) -- (pall);
  \draw[->] (pall) -- (sall);
\end{tikzpicture}

    \caption{Hierarchical summarization as a binary tree. Each leaf block $x_i$ is summarized once to produce $s_i=g(x_i)$. Internal nodes recursively summarize the concatenation of their children's stored strings, e.g., $s_{12}=g(s_1\concat s_2)$ and $s_{\mathrm{root}}=g(s_{1234}\concat s_{5678})$. The latent raw text under a node is denoted $S(u)$ for analysis, but the system ever only materializes the summaries $s_u$, making every merge auditable from bounded inputs.}
    \label{fig:hierarchical_summarization}
\end{figure}

Figure~\ref{fig:hierarchical_summarization} also locks in the notation we use later. Lowercase $x_i$ denotes the raw substrings carved out of the base document $x_0$; their summaries are $s_i=g(x_i)$. Every internal node carries two labels: $S(u)$ for the latent ``what text lives under here'' span and $s_u$ for the actual string we store and propagate. The recursion $s_u=g(s_{u_L}\concat s_{u_R})$ means parents ever only see $(s_{u_L},s_{u_R})$, so every merge fits inside the context window. The root ends up with $s_{\mathrm{root}}=Z^{(1)}(x)$ while $S(\text{root})=x$, making it easy to reason about fidelity from the leaves all the way up.

\paragraph{Remark (well-formed blocks).}
Throughout we assume the partition $(b_1,\ldots,b_k)$ is fixed in advance, consists of contiguous substrings of the underlying document, and respects whatever discourse boundaries the downstream task requires. In practice this means that each block is small enough to fit in context but large enough to carry coherent snippets of evidence (e.g., full sentences or paragraph segments). When we speak about ``realized'' leaves in the audit, we refer precisely to these blocks; all guarantees are conditioned on the scheduler never reordering or overlapping them. This is the only structural assumption on the data.

\subsection{Summarization Algebra}

If we were summarizing simple data structures, the requirements for $g$ would be straightforward. Consider the case where $\fstar$ is a count of keywords. Concatenating spans simply adds their counts: $\fstar(u \concat v) = \fstar(u) + \fstar(v)$. To preserve this information, a summarizer $g$ (acting as a ``sketch'') would need to support a merge operation that mimics addition. This is the essence of \emph{mergeable summaries} in database theory \citep{Agarwal2013MergeableTODS}:
\[
\fstar(A \cup B) \approx \text{Query}\big(\text{Merge}(S_A, S_B)\big),
\]
where $S_A$ denotes the stored sketch of set $A$ and $\text{Merge}(S_A,S_B)$ produces a sketch of $A\cup B$. For complex semantic tasks---event extraction, narrative summarization, or legal analysis---the merge operation is concatenation, which is non-commutative ($u \concat v \neq v \concat u$). Nevertheless, the algebraic intuition survives. We view $g$ as a \emph{summarization algebra} in which the merge operator is implemented by concatenation followed by another call to $g$. Ideally, we would have
\[
\metric\Big(\fstar(u \concat v), \;\fstar\big(g(u) \concat g(v)\big)\Big) = 0
\quad\text{for all }u,v\in\Strings,
\]
so that every hierarchy, schedule, or reduction behaves as a homomorphism with respect to $\fstar$. If this global property held, hierarchical summarization would be trivially safe. However, Large Language Models are not naturally associative algebraic operators. They are stochastic, sensitive to phrasing, and prone to ``drifting'' when context is segmented. We cannot assume this global structure exists; we must \emph{enforce} it via local checks that approximate the desired algebra.

\subsection{Consistency Conditions: Sufficiency, Merge Consistency, Idempotence}

Since we cannot verify the global property on massive corpora (which would require running $\fstar$ on the full text, defeating the purpose of summarization), we focus on properties that are \emph{local} and \emph{auditable}. We refer to them simply as the sufficiency, merge-consistency, and idempotence conditions. They act as structural invariants in theory and as prompt-engineering targets in practice.

\begin{convention}[Standing expectation convention]
Expectations $\E[\cdot]$ at a node $u$ are conditional on the realized subtrees below it. They average over the coin flips used to generate the summaries at $u$ and its children. All almost-sure statements are with respect to the same randomness.
\end{convention}

\paragraph{Condition 1: Sufficiency.}
The summary of every realized leaf must be a valid proxy for the raw block. Using the oracle equivalence notation, the theoretical requirement is
\begin{equation}\tag{C1}\label{eq:leaf_sufficiency}
\forall \text{ realized leaf } b:\quad g(b)\sim b.
\end{equation}
Once Condition~\eqref{eq:leaf_sufficiency} holds, a leaf summary may replace its source block anywhere because the two strings are equivalent at the oracle. In particular, if a block contains no evidence, both $b$ and $g(b)$ map to the neutral element $\emptyset$ in $\Y$; if a block introduces specific actors or labels, those items must survive in $g(b)$. If an LLM cannot satisfy this condition even on a single paragraph (e.g., by dropping a name), the hierarchy fails immediately. Section~\ref{sec:audit} therefore estimates the \emph{Empirical Violation Rate} of this condition by sampling leaves and checking how often the oracle disagrees.

\paragraph{Condition 2: Merge Consistency.}
The summarization pipeline must commute with concatenation modulo the oracle. To ensure the hierarchy structure does not corrupt the task signal, we require the following chain of equivalences for any $u,v\in\Strings$:
\begin{equation}\tag{C2}\label{eq:leaf_compatibility}
\forall u,v\in\Strings:\quad
\underbrace{u\concat v}_{\text{Raw}} \;\sim\;
\underbrace{g(u\concat v)}_{\text{Joint}} \;\sim\;
\underbrace{g\big(g(u)\concat g(v)\big)}_{\text{Disjoint}}.
\end{equation}
The first link ($u\concat v \sim g(u\concat v)$) asserts that jointly summarizing a span preserves its task value. The second link ($g(u\concat v) \sim g(g(u)\concat g(v))$) asserts that splitting the inputs before summarization is information-equivalent to processing them together. Because Condition~\eqref{eq:leaf_sufficiency} immediately yields $g(z)\sim z$ for any realized summary $z$, enforcing Condition~\eqref{eq:leaf_compatibility} on raw strings automatically controls every on-range pair. Condition~\eqref{eq:leaf_compatibility} is universal---it quantifies over every pair regardless of where it appears in the tree—but operationally we audit whichever link fits in context (Section~\ref{sec:audit}). Figure~\ref{fig:local_compatibility} shows how the two admissible routes align at a realized merge.

\begin{figure}[t]
    \centering
    \begin{tikzpicture}[node distance=8.5mm and 10.5mm,scale=0.95, every node/.style={transform shape}]
% Route 1 (Reference)
\node[doc] (a1) at (0,0) {$v$};
\node[doc] (b1) at (1.2,0) {$w$};
\node[plus] (plus1) at (0.6,1.1) {$\concat$};
\node[sum] (gab) at (0.6,2.3) {$g(v \concat w)$};
\node[op] (f1) at (0.6,3.5) {$\fstar(g(v \concat w))$};
\draw[->] (a1) -- node[left] {\scriptsize Raw} (plus1);
\draw[->] (b1) -- node[right] {\scriptsize Raw} (plus1);
\draw[->] (plus1) -- node[right] {\scriptsize Summarize} (gab);
\draw[->] (gab) -- node[right] {\scriptsize Oracle} (f1);

% Route 2 (Approximation)
\node[doc] (a2) at (4.6,0) {$v$};
\node[doc] (b2) at (5.8,0) {$w$};
\node[sum] (ga) at (4.6,1.1) {$g(v)$};
\node[sum] (gb) at (5.8,1.1) {$g(w)$};
\node[plus] (plus2) at (5.2,2.3) {$\concat$};
\node[sum] (ggab) at (5.2,3.5) {$g(g(v) \concat g(w))$};
\node[op] (f2) at (5.2,4.7) {$\fstar(g(g(v) \concat g(w)))$};
\draw[->] (a2) -- node[left] {\scriptsize Summarize} (ga);
\draw[->] (b2) -- node[right] {\scriptsize Summarize} (gb);
\draw[->] (ga) -- node[left] {\scriptsize Raw concat} (plus2);
\draw[->] (gb) -- node[right] {\scriptsize Raw concat} (plus2);
\draw[->] (plus2) -- node[right] {\scriptsize Summarize} (ggab);
\draw[->] (ggab) -- node[right] {\scriptsize Oracle} (f2);

\node[above=0.2cm of f1] {\scriptsize Reference path};
\node[above=0.2cm of f2] {\scriptsize Approximation path};

\draw[<->, dashed, thick] (f1) -- node[above, sloped] {\tiny $v\concat w \sim g(v\concat w) \sim g(g(v)\concat g(w))$} (f2);
    \end{tikzpicture}
    \caption{Merge consistency at a realized merge. The \emph{reference path} (left) concatenates the raw spans and summarizes once before applying the oracle. The \emph{approximation path} (right) summarizes each child, concatenates the summaries, and summarizes again before invoking the oracle. Condition~\eqref{eq:leaf_compatibility} enforces commutativity between concatenation and summarization, so both paths return oracle-equivalent values and the dashed equality is zero in expectation.}
    \label{fig:local_compatibility}
\end{figure}

\begin{remark}[Edge-wise compatibility]\label{rem:edgewise_compatibility}
Audits often instantiate Condition~\eqref{eq:leaf_compatibility} at a single realized edge $(v,w)$ by comparing the reference and pipeline routes explicitly (see the leaf-boundary case in Section~\ref{sec:audit}):
\begin{equation}\tag{E1}\label{eq:edgewise_equivalence}
\metric\!\big(\fstar(v\concat w),\,\fstar(g(g(v)\concat g(w)))\big)=0.
\end{equation}
Equation~\eqref{eq:edgewise_equivalence} probes the \emph{right} link of the merge-consistency chain using quantities that fit inside the context window. It complements the \emph{left} link ($v\concat w \sim g(v\concat w)$), which is testable whenever the raw concatenation is small enough to run directly through $g$.
\end{remark}

\paragraph{Condition 3: Idempotence.}
Hierarchical pipelines often include optional ``clean-up'' passes that re-summarize a stored output to smooth the style or reduce length. For this to be safe, the summarizer must be a fixed point on its own range:
\begin{equation}\tag{C3}\label{eq:leaf_idempotence}
\forall \text{ summary } s\in\operatorname{range}(g):\quad g(s)\sim s.
\end{equation}
Without Condition~\eqref{eq:leaf_idempotence}, repeatedly normalizing a string could cause the task information to drift. When it holds, any extra style edit, deduplication sweep, or DSPy controller pass is provably inert at the oracle.

Appendix~\ref{app:stability} constructs an explicit hierarchy where Conditions~(C1) and (C2) hold but $g$ drifts the label under repeated summarization; Condition~\eqref{eq:leaf_idempotence} rules out exactly this failure mode.

\subsection{Consequences: Processing Invariance}

The utility of these laws is that they allow us to transport trust from the leaves to the root. We do not need to verify the final summary against the whole book; we only need to verify that the ``assembly line'' is functioning correctly at each station.

\begin{theorem}[Inductive Preservation]\label{thm:one-pass}
Fix a document partition $x = b_1 \concat \dots \concat b_k$ and any full binary merge tree $T$. If \ref{eq:leaf_sufficiency} holds on every realized leaf and \ref{eq:leaf_compatibility} holds at every realized internal node, then the final summary preserves the oracle in expectation:\footnote{All expectation-based bounds in this paper assume the oracle space $(Y,\metric)$ has bounded diameter, i.e., $\sup_{y,y'}\metric(y,y')<\infty$. This is automatic for classification tasks (finite label sets) and bounded reward functions. The formal Lean development tracks this via \texttt{BoundedPseudoMetricSpace} typeclasses.}
\[
\E\left[\metric\big(\fstar(Z^{(1)}), \fstar(x)\big)\right] = 0.
\]
Furthermore, if \ref{eq:leaf_idempotence} holds, this preservation is maintained through any number of additional re-summarization rounds $Z^{(R)}$.
\end{theorem}

\noindent\textit{Proof:} See Appendix~\ref{app:proofs-core}, Proof~\ref{proof:thm-one-pass}.

\begin{corollary}[Schedule invariance]\label{cor:schedule}
For any fixed partition $(b_1,\ldots,b_k)$, every full binary tree on these leaves yields the same expected oracle whenever Conditions~\eqref{eq:leaf_sufficiency} and \eqref{eq:leaf_compatibility} hold on its realized edges. Balanced reductions and daisy chains are therefore interchangeable.
\end{corollary}

\noindent\textit{Proof:} See Appendix~\ref{app:proofs-core}, Proof~\ref{proof:cor-schedule}.

\begin{corollary}[Fold-of-folds invariance]\label{cor:folds}
Consider any two-level plan that first reduces contiguous ``folds'' and then reduces the fold summaries. If Conditions~\eqref{eq:leaf_sufficiency} and \eqref{eq:leaf_compatibility} hold on every realized edge and Condition~\eqref{eq:leaf_idempotence} holds on the intermediate summaries, the composite schedule preserves $\fstar$ in expectation regardless of the within-fold or over-fold parenthesizations.
\end{corollary}

\noindent\textit{Proof:} See Appendix~\ref{app:proofs-core}, Proof~\ref{proof:cor-folds}.

Together these results ensure that practitioners can pick whatever reduction schedule matches their infrastructure (balanced trees, daisy chains, fold-of-folds pipelines) without changing the expected oracle, provided the three conditions hold in the realized hierarchy and Condition~\eqref{eq:leaf_idempotence} keeps additional normalization rounds inert.

\begin{theorem}[Multi-round preservation]\label{thm:multi-round}
Assume Conditions~\eqref{eq:leaf_sufficiency} and \eqref{eq:leaf_compatibility} hold along the realized tree so that Theorem~\ref{thm:one-pass} yields $\E[\metric(\fstar(Z^{(1)}(x)),\fstar(x))]=0$. If, in addition, on-range stability \eqref{eq:leaf_idempotence} holds, then $\E[\metric(\fstar(Z^{(R)}(x)),\fstar(x))]=0$ for every $R\ge1$ with $Z^{(R+1)}(x):=g(Z^{(R)}(x))$.
\end{theorem}

\noindent\textit{Proof:} See Appendix~\ref{app:proofs-core}, Proof~\ref{proof:thm-multi-round}.

The theorem makes explicit why the audit's stability term matters: additional ``clean-up'' passes are provably inert only when Condition~\eqref{eq:leaf_idempotence} holds, so every estimate of the idempotence violation rate feeds directly into the global guarantee.

By satisfying these consistency conditions, we transform the summarizer $g$ into a globally consistent operator. The remaining challenge is operational: how do we verify these invariants hold without checking every single node?

\subsection{Safe Learning on Summaries}\label{sec:exist-learn-transport}

Assumption~\ref{ass:CF} ensures that downstream training signals depend on documents only through $\fstar(X)$. When a hierarchy satisfies these conditions, the oracle of each realized node matches the oracle of its source text, so any supervision computed on top of the hierarchy is statistically indistinguishable from supervision computed on full documents. Crucially, this information-equivalence \emph{only} holds under Assumption~\ref{ass:CF}; if rewards depend on surface style, summarization can distort the target regardless of the audit. Appendix~\ref{app:blackwell} formalizes the reduction via Blackwell sufficiency, and Appendix~\ref{sec:dpo} explains why we focus on Direct Preference Optimization (DPO) \citep{RafailovEtAl2023DPO}: DPO is a practical way to align models with researcher preferences, its Bradley--Terry backbone matches our oracle-based factorization, and its loss is smooth enough to incorporate the condition checks as optimization objectives. Once the audit certifies the tree (small $\hat{p}_{\text{suff}}$, $\hat{p}_{\text{assoc}}$, $\hat{p}_{\text{idem}}$), training on summaries achieves the same optimum as training on full text up to a gap controlled by the measured distortion, so the preference model inherits the guarantees proved above. Because $\fstar$ may encode arbitrary (but expensive) researcher judgments, this route lets us align to those judgments directly; aesthetic or stylistic goals outside $\fstar$ remain out of scope.

\subsection{Constructive Alignment via Local Bootstrapping}\label{sec:bootstrap-main}

These conditions are therefore more than diagnostic tools; they describe an optimization problem. We view the summarization system as a programmable policy $\mathcal{P}_\theta$ whose parameters $\theta$ collect prompt instructions, few-shot exemplars, or DSPy controller settings. Because Conditions~\eqref{eq:leaf_sufficiency} and \eqref{eq:leaf_compatibility} evaluate only the strings already stored at child nodes, each audited edge yields an inexpensive supervision signal: a success certifies $(\text{input},\text{summary})$ as a valid training pair, and a failure tells us exactly which constraint was violated.

In practice we maintain two light-weight modules. A \texttt{LeafSummarizer} maps raw blocks $b$ to $s_b=g(b)$ and is tuned to satisfy Condition~\eqref{eq:leaf_sufficiency}. A \texttt{MergeSummarizer} maps stored child summaries $(s_{u_L},s_{u_R})$ to $s_u$ and is tuned to satisfy Condition~\eqref{eq:leaf_compatibility}. A simple ``consistency bootstrap'' (an active-learning-style loop) alternates between them:
\begin{enumerate}[label=\roman*.]
    \item Sample leaves and runs of the merge tree. For each leaf, run \texttt{LeafSummarizer}, evaluate $\metric(\fstar(g(b)),\fstar(b))$, and, when necessary, re-generate or request a human rewrite until the oracle agrees. Certified pairs $(b,s_b)$ become part of the prompt context or fine-tuning buffer for the leaf module.
    \item Using the improved leaf module, generate child summaries $(s_{u_L},s_{u_R})$, apply \texttt{MergeSummarizer}, check $\metric(\fstar(g(s_{u_L}\concat s_{u_R})),\fstar(s_{u_L}\concat s_{u_R}))$, and add successful triples to the merge module's training set. Because the check fits in context, we can perform this loop offline or online without touching the raw subtree.
    \item Continue sampling until the empirical failure rates estimated by the audit fall below the tolerance required for deployment. Additional human examples can be injected at any time; they are especially useful on structures where repeated regeneration never satisfies the oracle.
\end{enumerate}
Appendix~\ref{app:bootstrap} provides a DSPy-oriented template for this loop. The main text only needs the structural fact: the same local comparisons that certify past runs also act as objective functions for the optimizer, allowing practitioners to tune $\theta$ with no extra labels beyond the oracle already assumed in our framework.
